\section{Machine-checked development}
\label{sec:formalization}

The repository includes a Lean~4 formalisation
\cite{deMouraUllrich2021Lean4} of the finite algebra and theorem composition
used in this paper. The \path{AspisFormal/} project kernel-checks integer value
conservation, relation closure from the gate residuals, the manifest-bound
Johnson/MCA regime, the complete finite soundness ledger, the
work-normalised endpoint, circle-generator order and fibre distinctness,
distribution-level masking lemmas, and the Poseidon2 known-answer bindings.
The theft-resistance module proves the final implication from its explicitly
named extractor and nullifier-binding interfaces. Every audited theorem uses
only \texttt{propext}, \texttt{Classical.choice}, and \texttt{Quot.sound}; the
development contains no \texttt{sorry} or custom axiom.

The proof-status table in \path{AspisFormal/README.md} separates kernel
theorems from cited interfaces. In particular, the published
Johnson/MCA/PCS, Fiat--Shamir, extractor, and simulation-extractability
results are imported at their mathematical statements rather than
reconstructed. Hash security and random-oracle modelling remain
cryptographic assumptions. This is the same boundary used by the paper's
security reductions.

After the q18/g37 mainnet execution, the project extended this architecture to
the V5 Tag-67 deployment candidate with Charon and Aeneas
\cite{HoProtzenko2022Aeneas}. The pipeline is
\[
  \text{production Rust}
  \longrightarrow
  \text{Aeneas-generated Lean}
  \longrightarrow
  \text{maintained protocol model}
\]
\[
  \longrightarrow\ \text{integration theorem}.
\]
The final V5 theorem composes the source-extracted Component-A execution at
the selected release schedule, general Component-B correspondence, the
actual-current Component-C four-round evaluator and public packed output, and
the Tag-67 work-wire/six-step verifier. The runtime recomputes both
GoodA and GoodB for every selected branch; a universal all-schedule
source-authentic Component-A theorem remains separate. The theorem's sole Tag-67
implementation-to-model premise identifies the stored hash-function
application with the maintained function on
\(\texttt{DOM\_GRIND}\mathbin\Vert\texttt{nonce\_le64}\). Wire parsing,
projection, the leading-zero predicate, and the ordered six checks are proved
inside the theorem.

The q18/g37 transaction and V5 candidate remain distinct artefacts. The former
is the mainnet execution studied in this paper. The V5 SBF later
finalized the atomic Tag-67 path on devnet at 1,335,952 CU, demonstrating the
stronger source-authentic verification architecture now present on the
repository's default branch. At this manuscript cutoff, V5 had not yet been
executed on mainnet.

\paragraph{Repository outcome addendum (25 July 2026).}
Subsequent to the manuscript-scoped q18/g37 result and the V5 freeze described
above, the exact V5 SBF (SHA-256
\nolinkurl{4cf3c1d5ddd47efa68875c0070247e007083c5c9bb2d5988db0d644a609edf40})
was deployed on mainnet-beta at slot \(434999519\).  Atomic Tag~67 finalized
at slot \(435019536\) with canonical nullifier-PDA bump \(255\); exact
signed-wire simulation and landed metadata both reported \(1{,}334{,}452\)
CU.  Its transaction signature is
\nolinkurl{EJviPgF12i9iK2CveVaQSMeFQqDMFPQ1iPRUYEwNQE3zGquTUZNJXPZEENorcQtsnQj1orFmH1TPsgdbR3vJ2fE}.
The corresponding proof and statement SHA-256 values are
\nolinkurl{330414df587974684643a6062d092db0519d746f0c7efe4ed2108775b685feaf}
and
\nolinkurl{0cdc34bc7f835640cff76d1085df9ba966df9f39eb228f3002f927cf30958113}.
This later chain outcome is recorded separately under
\path{release/aspis-v5-tag67-mainnet-v1/}; it does not change the q18/g37
parameters, reductions, or principal empirical claim studied in this paper.
